\chapter{GMRES and FGMRES}
\label{ch:gmres}

\chapterorient{The previous two chapters built the machinery of Krylov subspaces
and spent it on the friendliest possible case: a symmetric, positive-definite
operator, where conjugate gradients rides a three-term recurrence to the answer
on a fixed budget of storage. This chapter asks what happens when the operator is
\emph{not} symmetric---the driven analysis, the eigenmode inner solves, the
complex frequency-domain systems---where the short recurrence simply does not
exist. The answer is GMRES: keep the \emph{whole} Krylov basis, and at every step
pick the iterate that makes the residual as small as it can be. We build it from
the Arnoldi process up, watch a big minimization shrink to a tiny least-squares
problem, and see why the residual norm falls out for free at every step. Then we
pay off a promise made back in \cref{ch:operators}: when the preconditioner is
allowed to \emph{change every step}, plain GMRES breaks, and FGMRES---the same
algorithm carrying one extra basis---is exactly the repair the decision rule
predicts.}

\section{Why the short recurrence is gone}

Conjugate gradients (\cref{ch:cg}) is a small miracle of bookkeeping. To minimize
the energy of the error over the growing Krylov subspace $\mathcal{K}_m(\mat{A},
\vr_0)$, it never stores more than a handful of vectors: the new search direction
is $\mat{A}$-conjugate to all the previous ones, yet computing it requires only
the \emph{last} direction, because symmetry of $\mat{A}$ forces all the earlier
conjugacy conditions to hold automatically. The cost per step is fixed; the
storage is fixed; the algorithm is, in a precise sense, optimal on a budget.

That miracle has a single load-bearing hypothesis, and it is \emph{symmetry}. The
three-term recurrence at the heart of CG is the Lanczos recurrence, and Lanczos
tridiagonalizes $\mat{A}$ only because $\mat{A}=\mat{A}^{\top}$. The moment we drop
symmetry---and most of the operators in this book are not symmetric: the
frequency-domain system $\mat{A}(\omega)=\mat{K}+i\omega\mat{C}-\omega^2\mat{M}$ of
the driven analysis (\cref{ch:targets}) is complex and non-Hermitian; the shifted
operators inside an eigenvalue solve are indefinite---the recurrence collapses.
There is a theorem (Faber--Manteuffel) that makes this brutal: for a general
nonsymmetric $\mat{A}$, \emph{no} short recurrence can generate the optimal Krylov
iterate. You cannot have both optimality and bounded storage. You must give up one.

GMRES gives up bounded storage and keeps optimality. It commits to remembering the
\emph{entire} orthonormal basis of the Krylov subspace, and in exchange it delivers,
at every step, the genuinely best iterate the subspace can offer. The chapter is the
story of how it affords this---and, when the cost grows too large, how restarting
trades a controlled amount of optimality back for a memory bound.

\begin{notationbox}
Throughout, $\mat{A}$ is a general (possibly nonsymmetric, possibly complex
non-Hermitian) invertible operator on $\Real^N$ or $\Complex^N$; we solve
$\mat{A}\vx=\vb$. We write $\vx_0$ for the initial guess, $\vr_0=\vb-\mat{A}\vx_0$
for the initial residual, $\beta=\norm{\vr_0}$ for its norm, and
$\mathcal{K}_m=\mathcal{K}_m(\mat{A},\vr_0)=\operatorname{span}\{\vr_0,\mat{A}\vr_0,
\dots,\mat{A}^{m-1}\vr_0\}$ for the $m$-th Krylov subspace. The norm $\norm{\cdot}$
is always the Euclidean ($2$-)norm. Vectors $\vv_1,\dots,\vv_m$ are the Arnoldi
basis; $\vec e_1$ is the first coordinate vector.
\end{notationbox}

\subsection{The minimum-residual idea}

CG minimized the $\mat{A}$-norm of the error, a quantity that only makes sense---and
is only non-negative---when $\mat{A}$ is positive definite. For a general operator
that functional is meaningless. So GMRES minimizes something that always makes
sense: the ordinary Euclidean norm of the \emph{residual}.

\begin{keyidea}
The \textbf{minimum-residual principle.} At step $m$, GMRES chooses the iterate
$\vx_m$ from the affine Krylov subspace $\vx_0+\mathcal{K}_m$ that makes the
residual as small as possible:
\[
  \vx_m \;=\; \operatorname*{arg\,min}_{\vx\,\in\,\vx_0+\mathcal{K}_m}\ \norm{\vb-\mat{A}\vx}.
\]
Because $\mathcal{K}_m\subseteq\mathcal{K}_{m+1}$, the feasible set only grows from
one step to the next, so the minimum can only \emph{decrease}. The residual norms
are therefore \emph{monotonically non-increasing}: $\norm{\vr_0}\ge\norm{\vr_1}\ge
\norm{\vr_2}\ge\cdots$. GMRES never takes a step backward---a guarantee CG cannot
make for indefinite systems and that BiCGStab cannot make at all.
\end{keyidea}

\Cref{fig:minres} draws the principle. The affine subspace $\vx_0+\mathcal{K}_m$ is
a flat slice of $\Real^N$ that grows one dimension wider at each step. Over each
slice we minimize the residual; the minimizer is the point of the slice closest
(in the $\mat{A}$-image) to $\vb$. As the slice widens, the achievable minimum
drops, and the residual curve marches monotonically down.

\begin{figure}[t]
\centering
\begin{tikzpicture}[scale=1.0]
  % the ambient space hint
  \node[font=\scriptsize\itshape, simGray] at (-3.1,2.3) {$\Real^N$};
  % growing affine subspaces as nested line segments through x0
  \coordinate (x0) at (-2.6,-1.2);
  \fill[simInk] (x0) circle (1.4pt);
  \node[font=\scriptsize] at (x0) [below left=0pt and -1pt] {$\vx_0$};
  % K_1 : a short segment
  \draw[simGreen, thick] (x0) -- ++(2.2,0.5)
    node[pos=1.0, right, font=\scriptsize, simGreen] {$\vx_0+\mathcal{K}_1$};
  % K_2 : a wider wedge (a 2D slice drawn as a parallelogram)
  \draw[simBlue, thick, fill=simBgBlue, fill opacity=0.45]
    (x0) -- ++(2.2,0.5) -- ++(0.2,1.7) -- ++(-2.2,-0.5) -- cycle;
  \node[font=\scriptsize, simBlue] at (0.2,1.0) {$\vx_0+\mathcal{K}_2$};
  % the iterates, each the closest point in its slice
  \fill[simGreen] (-0.85,-0.78) circle (1.4pt);
  \node[font=\scriptsize, simGreen] at (-0.85,-0.78) [below=1pt] {$\vx_1$};
  \fill[simRust] (-0.2,0.35) circle (1.6pt);
  \node[font=\scriptsize, simRust] at (-0.2,0.35) [right=1pt] {$\vx_2$};
  % the target (in A-image, schematically): the "true solution" direction
  \node[draw=simGold, fill=simBgGold, thick, rounded corners=2pt,
        font=\scriptsize, inner sep=2pt] (sol) at (2.9,1.9) {$\vx_\star=\mat{A}^{-1}\vb$};
  % residual arrows: from each iterate "toward" the solution, shrinking
  \draw[backflow, shorten >=2pt] (-0.85,-0.78) -- (sol)
    node[pos=0.45, above, font=\scriptsize, simRust, sloped] {$\norm{\vr_1}$};
  \draw[flow, simRust, shorten >=2pt] (-0.2,0.35) -- (sol)
    node[pos=0.5, below, font=\scriptsize, simRust, sloped] {$\norm{\vr_2}<\norm{\vr_1}$};
\end{tikzpicture}
\caption[The minimum-residual idea]{The minimum-residual principle. At each step
the feasible set is the affine Krylov subspace $\vx_0+\mathcal{K}_m$---a flat slice
of $\Real^N$ that gains one dimension per step (here $\mathcal{K}_1\subset
\mathcal{K}_2$). GMRES returns the point $\vx_m$ of that slice whose residual
$\norm{\vb-\mat{A}\vx_m}$ is smallest. Because each slice contains the previous one,
the achievable minimum can only shrink: $\norm{\vr_2}\le\norm{\vr_1}\le\norm{\vr_0}$.
The residual decreases monotonically by construction---GMRES never moves backward.}
\label{fig:minres}
\end{figure}

There is one immediate problem, and the rest of the chapter is its solution. The
minimization is posed over a subspace of $\Real^N$, and $N$ is enormous---millions
of unknowns for a real mesh. We cannot solve an $N$-dimensional least-squares
problem at every step. The genius of GMRES is that we never have to: by choosing
the \emph{right basis} for $\mathcal{K}_m$, the giant minimization collapses to a
tiny one of size $m+1$, where $m$ is the step count, a few dozen at most. That
right basis is built by the Arnoldi process.

\section{Arnoldi: an orthonormal basis for the Krylov subspace}

The raw Krylov vectors $\vr_0,\mat{A}\vr_0,\mat{A}^2\vr_0,\dots$ span
$\mathcal{K}_m$ but are a terrible basis: they line up more and more with the
dominant eigenvector of $\mat{A}$, becoming numerically parallel within a handful
of steps. We need an \emph{orthonormal} basis instead, and we need to build it one
vector at a time, because we discover $\mathcal{K}_m$ one matvec at a time. The
Arnoldi process does exactly this: it is Gram--Schmidt orthogonalization (the full
treatment is \cref{ch:orthog}) applied to the Krylov sequence as it is generated.

\subsection{The process}

Start with the normalized initial residual $\vv_1=\vr_0/\beta$. At each step, take
the most recently produced basis vector $\vv_j$, hit it with $\mat{A}$ to push into
the next Krylov dimension, then \emph{orthogonalize} the result against every basis
vector built so far and normalize what remains:

\begin{lstlisting}[style=l4style]
-- one Arnoldi step: extend an orthonormal Krylov basis by one vector
arnoldi_step :: LinOp -> Vec[] -> Int -> (Vec, Float[], Float)
arnoldi_step A V j =
  let w0 = matvec A (V[j])           -- push into the next Krylov dimension
  let (w, h) = orthogonalize w0 V    -- subtract off components along V[0..j]
  --   h[k] = <w0, V[k]>  are the Arnoldi coefficients (the Hessenberg column)
  let hjj1 = norm w                  -- the new sub-diagonal entry  ||w||
  let v_next = w `scale` (1 / hjj1)  -- normalize: the new basis vector V[j+1]
  in (v_next, h, hjj1)
\end{lstlisting}

\noindent The coefficients $h_{kj}=\inner{\mat{A}\vv_j}{\vv_k}$ for $k\le j$ record
how much of $\mat{A}\vv_j$ pointed along each existing basis vector; the
sub-diagonal entry $h_{j+1,j}=\norm{\vw}$ records the size of the genuinely new
direction. Collect these numbers and a remarkable structure appears.

\begin{figure}[t]
\centering
\begin{tikzpicture}[scale=0.92]
  % ---- LEFT: the tall-skinny basis V_{m+1} ----
  \begin{scope}
    \node[font=\bfseries\small, simInk] at (0,3.05) {The basis $\mat{V}_{m+1}$};
    \draw[simBlue, thick, fill=simBgBlue] (-1.0,-2.6) rectangle (1.0,2.6);
    % columns
    \foreach \x in {-0.6,-0.2,0.2,0.6}{
      \draw[simBlue!55] (\x,-2.6) -- (\x,2.6);}
    \node[font=\scriptsize, simGray, rotate=90] at (0,0) {$N$ rows (huge)};
    \node[font=\scriptsize, simBlue] at (0,2.85) {$m{+}1$ columns};
    \node[font=\scriptsize\itshape, simGray, align=center] at (0,-3.0)
      {orthonormal\\columns $\vv_1\dots\vv_{m+1}$};
  \end{scope}
  % ---- the Arnoldi relation glyph ----
  \node[font=\Large] at (2.0,0) {$\mat{A}\,\mat{V}_m=$};
  % ---- RIGHT: V_{m+1} times the small Hessenberg ----
  \begin{scope}[shift={(4.5,0)}]
    % small V_{m+1} again
    \draw[simBlue, thick, fill=simBgBlue] (-0.8,-2.0) rectangle (0.5,2.0);
    \node[font=\scriptsize, simBlue] at (-0.15,2.25) {$\mat{V}_{m+1}$};
    \node[font=\large] at (1.0,0) {$\cdot$};
  \end{scope}
  % ---- the Hessenberg H-bar (small, (m+1) x m) ----
  \begin{scope}[shift={(6.6,0)}]
    \node[font=\bfseries\small, simInk] at (0,3.05) {$\bar{\mat{H}}_m$ (small)};
    \draw[simRust, thick, fill=simBgRust] (-0.95,-1.0) rectangle (0.95,2.0);
    \node[font=\scriptsize, simRust] at (0,2.25) {$m$ cols};
    \node[font=\scriptsize, simRust, rotate=90] at (-1.25,0.5) {$m{+}1$ rows};
    % shade the upper-Hessenberg pattern: filled on/above first subdiagonal
    \fill[simRust!30] (-0.95,2.0) -- (0.95,2.0) -- (0.95,-1.0) -- (-0.55,-1.0)
      -- (-0.95,-0.4) -- cycle;
    % subdiagonal staircase
    \draw[simRust!75, very thick] (-0.95,-0.4) -- (-0.55,-1.0);
    \node[font=\scriptsize\itshape, simGray, align=center] at (0,-1.7)
      {upper-Hessenberg:\\one nonzero\\sub-diagonal};
  \end{scope}
\end{tikzpicture}
\caption[Arnoldi builds a tall basis and a small Hessenberg]{The Arnoldi process
produces two objects of utterly different size. \emph{Left:} the orthonormal basis
$\mat{V}_{m+1}$, an $N\times(m{+}1)$ matrix---tall and skinny, $N$ in the millions
but only $m{+}1$ columns. \emph{Right:} the Arnoldi relation
$\mat{A}\mat{V}_m=\mat{V}_{m+1}\bar{\mat{H}}_m$ expresses the action of $\mat{A}$ on
the basis entirely through the \emph{small} $(m{+}1)\times m$ upper-Hessenberg
matrix $\bar{\mat{H}}_m$, whose only nonzeros are on and above the first
sub-diagonal. The huge operator $\mat{A}$ is, within the Krylov subspace, captured
by a matrix a few dozen entries on a side. That compression is what makes the
minimization tractable.}
\label{fig:arnoldi}
\end{figure}

\subsection{The Arnoldi relation}

Read the orthogonalization step backward. We computed $\mat{A}\vv_j$, then wrote it
as its components along the existing basis plus a new normalized direction:
\[
  \mat{A}\vv_j \;=\; \sum_{k=1}^{j} h_{kj}\,\vv_k \;+\; h_{j+1,j}\,\vv_{j+1}.
\]
This says $\mat{A}\vv_j$ is a combination of $\vv_1,\dots,\vv_{j+1}$ only---the
matvec never reaches past the next basis vector. Stacking these identities for
$j=1,\dots,m$ into matrix form gives the \emph{Arnoldi relation}, the single most
important equation in the chapter:

\begin{theorem}[Arnoldi relation]
\label{thm:arnoldi}
Let $\mat{V}_m=[\vv_1\ \cdots\ \vv_m]$ be the $N\times m$ matrix of Arnoldi basis
vectors and $\mat{V}_{m+1}=[\mat{V}_m\ \vv_{m+1}]$ the same with one extra column.
Let $\bar{\mat{H}}_m$ be the $(m{+}1)\times m$ matrix whose entry in row $k$,
column $j$ is the Arnoldi coefficient $h_{kj}$ (and zero where $k>j+1$). Then
\[
  \mat{A}\,\mat{V}_m \;=\; \mat{V}_{m+1}\,\bar{\mat{H}}_m,
\]
and $\bar{\mat{H}}_m$ is \emph{upper-Hessenberg}: $h_{kj}=0$ whenever $k>j+1$, so
its only sub-diagonal is the first one.
\end{theorem}

The Hessenberg structure is not a coincidence; it \emph{is} the statement that
$\mat{A}\vv_j$ reaches only as far as $\vv_{j+1}$. (When $\mat{A}$ is symmetric,
$\bar{\mat{H}}_m$ becomes tridiagonal---one sub-diagonal \emph{and} one
super-diagonal---and the orthogonalization collapses to the three-term Lanczos
recurrence of \cref{ch:cg}. Hessenberg is the nonsymmetric generalization;
tridiagonal is the symmetric special case.) \Cref{fig:arnoldi} contrasts the two
sizes: $\mat{V}_{m+1}$ is enormous in one dimension, but $\bar{\mat{H}}_m$ is tiny
in both.

\subsection{From a big minimization to a small one}

Now we cash in. Any candidate iterate in the affine subspace can be written
$\vx=\vx_0+\mat{V}_m\vy$ for some coordinate vector $\vy\in\Real^m$---the columns of
$\mat{V}_m$ span $\mathcal{K}_m$, so $\mat{V}_m\vy$ ranges over all of it as $\vy$
ranges over $\Real^m$. Substitute and compute the residual, using $\vr_0=\beta\vv_1$
and the Arnoldi relation:
\begin{align*}
  \vb-\mat{A}\vx
    &= \vr_0 - \mat{A}\mat{V}_m\vy
     = \beta\vv_1 - \mat{V}_{m+1}\bar{\mat{H}}_m\vy \\
    &= \mat{V}_{m+1}\bigl(\beta\vec e_1 - \bar{\mat{H}}_m\vy\bigr),
\end{align*}
where the last step pulls out $\mat{V}_{m+1}$ and uses $\vv_1=\mat{V}_{m+1}\vec e_1$.
Because the columns of $\mat{V}_{m+1}$ are orthonormal, multiplying by
$\mat{V}_{m+1}$ preserves the Euclidean norm. Therefore
\[
  \norm{\vb-\mat{A}\vx}
    \;=\; \norm{\mat{V}_{m+1}\bigl(\beta\vec e_1-\bar{\mat{H}}_m\vy\bigr)}
    \;=\; \norm{\beta\vec e_1-\bar{\mat{H}}_m\vy}.
\]

\begin{keyidea}
The minimum-residual problem in $\Real^N$ collapses to a least-squares problem in
$\Real^{m+1}$:
\[
  \min_{\vx\,\in\,\vx_0+\mathcal{K}_m}\norm{\vb-\mat{A}\vx}
  \;=\;
  \min_{\vy\,\in\,\Real^m}\ \norm{\beta\,\vec e_1-\bar{\mat{H}}_m\,\vy}.
\]
The orthonormality of the Arnoldi basis turned a minimization over a subspace of
the \emph{huge} space into one over $\vy$ of dimension $m$---a problem whose data is
the \emph{tiny} $(m{+}1)\times m$ Hessenberg matrix $\bar{\mat{H}}_m$ and the scalar
$\beta$. We solve this small problem for $\vy_m$, and only then form the actual
iterate $\vx_m=\vx_0+\mat{V}_m\vy_m$ in the big space---once, at the end.
\end{keyidea}

\Cref{fig:reduction} draws the reduction as a funnel: the $N$-dimensional residual
minimization on the left, the orthonormal-basis change in the middle, the
$(m{+}1)$-dimensional least-squares problem on the right. Everything expensive about
$\mat{A}$ and $N$ has been squeezed into building $\mat{V}_m$ and $\bar{\mat{H}}_m$;
the actual optimization happens in a space small enough to hold in a cache line.

\begin{figure}[t]
\centering
\begin{tikzpicture}[node distance=8mm]
  % BIG problem (left)
  \node[draw=simBlue, fill=simBgBlue, thick, rounded corners=2pt,
        minimum width=34mm, minimum height=20mm, align=center, font=\footnotesize]
        (big) at (0,0)
    {\textbf{big minimization}\\[2pt]$\displaystyle\min_{\vx\in\vx_0+\mathcal{K}_m}$\\[1pt]
     $\norm{\vb-\mat{A}\vx}$\\[2pt]\itshape\scriptsize over a subspace of $\Real^N$};
  % the lever: orthonormal basis + Arnoldi relation
  \node[stage, minimum width=30mm, minimum height=16mm, right=12mm of big, align=center]
        (lever)
    {\textbf{Arnoldi}\\[1pt]\footnotesize $\vx=\vx_0+\mat{V}_m\vy$\\\footnotesize
     $\mat{A}\mat{V}_m=\mat{V}_{m+1}\bar{\mat{H}}_m$\\\footnotesize
     $\norm{\mat{V}_{m+1}\vu}=\norm{\vu}$};
  \draw[flow] (big) -- node[above,font=\scriptsize]{change of} node[below,font=\scriptsize]{basis} (lever);
  % SMALL problem (right)
  \node[draw=simRust, fill=simBgRust, thick, rounded corners=2pt,
        minimum width=34mm, minimum height=20mm, align=center, font=\footnotesize]
        (small) at ($(lever.east)+(2.0,0)$)
    {\textbf{small least-squares}\\[2pt]$\displaystyle\min_{\vy\in\Real^m}$\\[1pt]
     $\norm{\beta\vec e_1-\bar{\mat{H}}_m\vy}$\\[2pt]\itshape\scriptsize
     dimension $m{+}1$ (tiny)};
  \draw[flow] (lever) -- (small);
  % annotation
  \node[font=\scriptsize\itshape, simGray, align=center, below=4mm of lever, text width=34mm]
    {the same minimum value,\\two sizes of problem};
\end{tikzpicture}
\caption[The big-to-small reduction]{The engine of GMRES. The minimum-residual
problem on the left lives in $\Real^N$ with $N$ in the millions. Writing every
candidate as $\vx_0+\mat{V}_m\vy$ and invoking the Arnoldi relation
(\cref{thm:arnoldi}) plus the norm-preservation of the orthonormal basis
$\mat{V}_{m+1}$ turns it into the least-squares problem on the right, whose data is
just the $(m{+}1)\times m$ Hessenberg matrix and the scalar $\beta$. Both problems
have \emph{exactly the same minimum value}; only the second is one we can afford to
solve at every step.}
\label{fig:reduction}
\end{figure}

\begin{workedexample}{Two Arnoldi steps on a small nonsymmetric operator}{arnoldi}
Take the nonsymmetric $3\times3$ operator and right-hand side
\[
  \mat{A}=\begin{bmatrix} 2 & 1 & 0\\ 0 & 3 & 1\\ 1 & 0 & 4\end{bmatrix},
  \qquad
  \vb=\begin{bmatrix}1\\0\\0\end{bmatrix},
  \qquad \vx_0=\mathbf{0}.
\]
$\mat{A}$ is not symmetric ($a_{31}=1\ne0=a_{13}$), so CG does not apply. We run two
Arnoldi steps by hand.

\textbf{Setup.} The initial residual is $\vr_0=\vb-\mat{A}\vx_0=\vb=[1,0,0]^{\top}$,
so $\beta=\norm{\vr_0}=1$ and the first basis vector is
$\vv_1=\vr_0/\beta=[1,0,0]^{\top}$.

\textbf{Step 1.} Push and orthogonalize. The matvec is
\[
  \mat{A}\vv_1=\begin{bmatrix}2\\0\\1\end{bmatrix}.
\]
Project onto $\vv_1$: the coefficient is $h_{11}=\inner{\mat{A}\vv_1}{\vv_1}=2$.
Subtract it off:
\[
  \vw=\mat{A}\vv_1-h_{11}\vv_1=\begin{bmatrix}2\\0\\1\end{bmatrix}
     -2\begin{bmatrix}1\\0\\0\end{bmatrix}=\begin{bmatrix}0\\0\\1\end{bmatrix}.
\]
The new sub-diagonal entry is $h_{21}=\norm{\vw}=1$, and the new basis vector is
$\vv_2=\vw/h_{21}=[0,0,1]^{\top}$. (It is orthonormal to $\vv_1$, as it must be.)

\textbf{Step 2.} Push and orthogonalize again. The matvec is
\[
  \mat{A}\vv_2=\begin{bmatrix}0\\1\\4\end{bmatrix}.
\]
Project onto both existing basis vectors:
$h_{12}=\inner{\mat{A}\vv_2}{\vv_1}=0$ and
$h_{22}=\inner{\mat{A}\vv_2}{\vv_2}=4$. Subtract both:
\[
  \vw=\mat{A}\vv_2-h_{12}\vv_1-h_{22}\vv_2
     =\begin{bmatrix}0\\1\\4\end{bmatrix}-0-4\begin{bmatrix}0\\0\\1\end{bmatrix}
     =\begin{bmatrix}0\\1\\0\end{bmatrix}.
\]
The sub-diagonal entry is $h_{32}=\norm{\vw}=1$, and $\vv_3=[0,1,0]^{\top}$.

\textbf{Read off the Hessenberg.} Assembling the coefficients into the
$3\times2$ upper-Hessenberg matrix:
\[
  \mat{V}_2=\begin{bmatrix}1&0\\0&0\\0&1\end{bmatrix},
  \qquad
  \bar{\mat{H}}_2=\begin{bmatrix} h_{11}&h_{12}\\ h_{21}&h_{22}\\ 0&h_{32}\end{bmatrix}
   =\begin{bmatrix} 2 & 0\\ 1 & 4\\ 0 & 1\end{bmatrix}.
\]
Notice the shape: $\bar{\mat{H}}_2$ is $(m{+}1)\times m=3\times2$, with a single
nonzero sub-diagonal (the $1$'s in positions $(2,1)$ and $(3,2)$)---exactly the
upper-Hessenberg pattern of \cref{thm:arnoldi}. We have compressed the action of the
$3\times3$ operator on $\mathcal{K}_2$ into these six numbers. We solve the
least-squares problem they pose in \cref{wex:leastsq}.
\end{workedexample}

\section{The small least-squares problem and the free residual}

We are left with $\min_{\vy}\norm{\beta\vec e_1-\bar{\mat{H}}_m\vy}$, an
overdetermined least-squares problem ($m{+}1$ equations, $m$ unknowns). The
textbook solution is a QR factorization of $\bar{\mat{H}}_m$: write
$\bar{\mat{H}}_m=\mat{Q}\mat{R}$ with $\mat{Q}$ orthonormal and $\mat{R}$ upper
triangular, then the minimizer solves $\mat{R}\vy=\mat{Q}^{\top}(\beta\vec e_1)$ by
back-substitution. But re-factorizing from scratch at every step would waste the
fact that $\bar{\mat{H}}_m$ grows by exactly one column per step. We can do far
better, and the better way hands us the residual norm for free.

\subsection{The running QR by Givens rotations}

A Hessenberg matrix is already \emph{almost} triangular: its only obstacle to upper
triangularity is the single sub-diagonal. Each new column, when it arrives, adds one
sub-diagonal entry that must be annihilated. A \emph{Givens rotation}---a $2\times2$
plane rotation acting on two adjacent rows---annihilates exactly one entry while
leaving the triangular structure above it intact. So we maintain the QR factorization
\emph{incrementally}: when column $j$ arrives, we replay the rotations generated for
columns $1,\dots,j-1$ onto it (bringing it into the established triangular frame),
then generate one new rotation to kill its fresh sub-diagonal entry. The mechanics of
the rotations---how each $2\times2$ rotation is computed and applied---are the
business of \cref{ch:givens}; here we use only the result.

\begin{figure}[t]
\centering
\begin{tikzpicture}[scale=0.95]
  % the growing triangular factor R, with one new column arriving
  \node[font=\bfseries\small, simInk] at (0,2.55) {Running QR: one new column per step};
  % already-triangular block
  \fill[simBlue!18] (-2.4,-1.4) -- (0.2,-1.4) -- (0.2,1.6) -- (-2.4,1.6) -- cycle;
  \draw[simBlue, thick] (-2.4,-1.4) rectangle (0.2,1.6);
  % triangular staircase inside
  \draw[simBlue!70, thick] (-2.4,1.6) -- (0.2,1.6) -- (0.2,-1.0) -- (-2.0,1.2) -- cycle;
  \node[font=\scriptsize, simBlue] at (-1.3,0.5) {already};
  \node[font=\scriptsize, simBlue] at (-1.3,0.1) {triangular};
  % the new column (highlighted)
  \draw[simRust, very thick, fill=simBgRust] (0.2,-1.4) rectangle (0.55,1.6);
  \node[font=\scriptsize, simRust, rotate=90] at (0.38,0.1) {new column $j$};
  % the offending subdiagonal entry
  \fill[simRust] (0.38,-1.15) circle (2pt);
  \node[font=\scriptsize, simRust, align=left] at (1.9,-1.15)
    {one sub-diagonal\\entry to annihilate};
  \draw[refedge] (1.0,-1.15) -- (0.5,-1.15);
  % the rotation
  \draw[-{Stealth[length=2mm]}, simGreen, thick] (1.1,0.6) arc (-40:40:0.5);
  \node[font=\scriptsize, simGreen, align=left] at (2.2,0.6)
    {generate one\\Givens rotation};
  % the RHS column on the far right, with the residual entry glowing
  \draw[simGold, thick, fill=simBgGold] (3.6,-1.4) rectangle (3.95,1.6);
  \node[font=\scriptsize, simGold!70!black, rotate=90] at (3.78,0.55) {rotated RHS $s$};
  \fill[simGold!80!black] (3.78,-1.15) circle (2.4pt);
  \node[font=\scriptsize\bfseries, simGold!60!black, align=left] at (5.0,-1.15)
    {$|s_{m+1}|=\norm{\vr_m}$\\\itshape\mdseries the residual,\\\itshape\mdseries for free};
  \draw[refedge] (4.2,-1.15) -- (3.9,-1.15);
\end{tikzpicture}
\caption[The running QR and the free residual readout]{The incremental QR. Each
Arnoldi step appends one column to the Hessenberg matrix; the already-computed
Givens rotations are replayed onto it, then one new rotation annihilates its single
fresh sub-diagonal entry, extending the triangular factor by one column. The same
new rotation is applied to the rotated right-hand side $s$ (initialized to
$\beta\vec e_1$): because every rotation is an isometry, the magnitude of its last
entry, $|s_{m+1}|$, equals the least-squares residual norm $\norm{\vr_m}$
\emph{exactly}---available at every step \emph{without} forming $\vx_m$ or computing
a single matvec. The rotation mechanics are developed in \cref{ch:givens}.}
\label{fig:runningqr}
\end{figure}

\subsection{The residual is free}

The payoff is the reason GMRES is built this way. Alongside the Hessenberg matrix we
carry a small right-hand-side vector $s$, initialized to $\beta\vec e_1$, and we
apply each Givens rotation to $s$ as we generate it. Because every rotation is an
isometry (it preserves length), the transformed system has the same residual norm as
the original. And after triangularization, the least-squares residual is precisely
the magnitude of the \emph{last} entry of $s$:
\[
  \norm{\vr_m}\;=\;\min_{\vy}\norm{\beta\vec e_1-\bar{\mat{H}}_m\vy}\;=\;\abs{s_{m+1}}.
\]

\begin{keyidea}
The GMRES residual norm is available \textbf{for free}, every step, as a single
scalar. After the running QR processes column $m$, the magnitude $|s_{m+1}|$ of the
last entry of the rotated right-hand side \emph{equals} $\norm{\vr_m}$
exactly---with no matvec, no explicit residual vector, and without ever forming the
iterate $\vx_m$. The convergence test reads off $|s_{m+1}|$ and compares it to the
tolerance. Only when the test passes (or we restart) do we pay the one-time cost of
back-solving the small triangular system for $\vy_m$ and forming
$\vx_m=\vx_0+\mat{V}_m\vy_m$ in the big space.
\end{keyidea}

This is what the source actually does: Palace's GMRES carries rotation registers
\code{cs} and \code{sn} and a rotated right-hand side \code{s}, applies each new
rotation to \code{s}, and reads the convergence test directly off
$\beta=|s_{j+1}|$---no residual is ever evaluated explicitly. The back-solve that
turns $\vy$ into a correction $\mat{V}_m\vy$ fires only at convergence or at a
restart boundary. We trace the small problem end to end now.

\begin{workedexample}{Solving the $3\times2$ least-squares and reading the residual}{leastsq}
Continue \cref{wex:arnoldi}, with $\beta=1$ and
\[
  \bar{\mat{H}}_2=\begin{bmatrix} 2 & 0\\ 1 & 4\\ 0 & 1\end{bmatrix},
  \qquad
  \beta\vec e_1=\begin{bmatrix}1\\0\\0\end{bmatrix}.
\]
We solve $\min_{\vy}\norm{\beta\vec e_1-\bar{\mat{H}}_2\vy}$ by the running QR,
generating two Givens rotations $\mat{G}_1,\mat{G}_2$ and watching the residual
appear in the rotated right-hand side $s$.

\textbf{Rotation 1} acts on rows $1,2$ to kill the $(2,1)$ entry. With diagonal
$\rho=2$ and sub-diagonal $\sigma=1$, the rotation has
$c_1=\rho/\sqrt{\rho^2+\sigma^2}=2/\sqrt5$ and
$s_1=\sigma/\sqrt{\rho^2+\sigma^2}=1/\sqrt5$. Applying $\mat{G}_1$ to the first
column of $\bar{\mat{H}}_2$ turns $(2,1)$ into $\sqrt{2^2+1^2}=\sqrt5$ on the
diagonal and $0$ below. Applying $\mat{G}_1$ to the second column
$[0,4,1]^{\top}$ mixes rows $1,2$: the new $(1,2)$ entry is
$c_1\!\cdot\!0+s_1\!\cdot\!4=4/\sqrt5$ and the new $(2,2)$ entry is
$-s_1\!\cdot\!0+c_1\!\cdot\!4=8/\sqrt5$. Applying $\mat{G}_1$ to the right-hand side
$s=[1,0,0]^{\top}$ gives
$s\leftarrow[c_1,\,-s_1,\,0]^{\top}=[2/\sqrt5,\,-1/\sqrt5,\,0]^{\top}$. After
rotation 1 the working system is
\[
  \mat{R}^{(1)}=\begin{bmatrix} \sqrt5 & 4/\sqrt5\\ 0 & 8/\sqrt5\\ 0 & 1\end{bmatrix},
  \qquad
  s=\begin{bmatrix} 2/\sqrt5\\ -1/\sqrt5\\ 0\end{bmatrix}.
\]

\textbf{Rotation 2} acts on rows $2,3$ to kill the $(3,2)$ entry. The diagonal is
$\rho=8/\sqrt5$ and the sub-diagonal $\sigma=1$, so
$r=\sqrt{\rho^2+\sigma^2}=\sqrt{64/5+1}=\sqrt{69/5}=\sqrt{69}/\sqrt5$, giving
$c_2=\rho/r=8/\sqrt{69}$ and $s_2=\sigma/r=\sqrt5/\sqrt{69}$. Applying $\mat{G}_2$ to
the right-hand side mixes rows $2,3$: the new entry $2$ is
$c_2\!\cdot\!(-1/\sqrt5)+s_2\!\cdot\!0=-8/\sqrt{5\cdot69}$, and---the entry we want---
the new entry $3$ is
\[
  s_3 \;=\; -s_2\cdot\!\left(-\tfrac{1}{\sqrt5}\right)+c_2\cdot 0
        \;=\; \frac{\sqrt5}{\sqrt{69}}\cdot\frac{1}{\sqrt5}
        \;=\; \frac{1}{\sqrt{69}}.
\]

\textbf{Read the residual for free.} The residual norm after two steps is
\[
  \norm{\vr_2}\;=\;\abs{s_3}\;=\;\frac{1}{\sqrt{69}}\;\approx\;0.1204,
\]
read directly off the last entry of the rotated right-hand side---no matvec, no
explicit residual vector. (The starting residual was $\norm{\vr_0}=\beta=1$, so two
GMRES steps cut it by a factor of about $8.3$.)

\textbf{Recover the iterate, once.} Only now do we back-solve the upper-triangular
$2\times2$ system $\mat{R}\vy=[s_1,s_2]^{\top}$ for the coordinate vector $\vy$. The
upper $2\times2$ block of the rotated Hessenberg is $\mat{R}=\bigl[\begin{smallmatrix}
\sqrt5 & 4/\sqrt5\\ 0 & 8/\sqrt{69}\cdot\sqrt5/\sqrt5\end{smallmatrix}\bigr]$ (with
$(2,2)=r=\sqrt{69}/\sqrt5$ after rotation~2), and the right-hand side is
$[2/\sqrt5,\,-8/\sqrt{5\cdot69}]^{\top}$. Back-substitution gives
$y_2=(-8/\sqrt{5\cdot69})/(\sqrt{69}/\sqrt5)=-8/69$ and
$y_1=(2/\sqrt5-(4/\sqrt5)y_2)/\sqrt5=(2+32/69)/5=170/345=34/69$. The iterate is
\[
  \vx_2=\vx_0+\mat{V}_2\vy=\begin{bmatrix}1&0\\0&0\\0&1\end{bmatrix}
    \begin{bmatrix}34/69\\-8/69\end{bmatrix}=\begin{bmatrix}34/69\\0\\-8/69\end{bmatrix}.
\]
As a check, $\mat{A}\vx_2=[60/69,\,-8/69,\,2/69]^{\top}$ and the true residual
$\vb-\mat{A}\vx_2=[9/69,\,8/69,\,-2/69]^{\top}$ has norm
$\sqrt{81+64+4}/69=\sqrt{149}/69\approx0.1768$. This differs from $|s_3|$ because the
$3\times3$ Krylov space is not yet exhausted; in exact arithmetic GMRES reaches the
exact solution at step $m=3=N$, where the residual hits zero. The point stands: the
\emph{small} problem's residual $|s_3|$ was free, and the iterate was formed only once.
\end{workedexample}

\section{Restart: bounding the cost of remembering everything}

GMRES bought optimality by remembering the entire basis. That memory is not free,
and as $m$ grows the bill comes due on two fronts.

\begin{itemize}[nosep]
  \item \textbf{Storage grows linearly in $m$.} The basis $\mat{V}_m$ holds $m$
  vectors of length $N$, so storage is $O(mN)$. At step 100 of a million-unknown
  solve, that is a hundred million-entry vectors---hundreds of megabytes, just for
  the basis.
  \item \textbf{Work grows quadratically in $m$.} Orthogonalizing the new vector
  against all $m$ existing ones costs $O(mN)$ per step, so the cumulative
  orthogonalization work through step $m$ is $O(m^2N)$. The matvecs are only $O(mN)$;
  it is the \emph{orthogonalization against an ever-growing basis} that dominates.
\end{itemize}

Both costs are unbounded in $m$, which is intolerable for a solve that might need
hundreds of iterations. The standard remedy is to cap $m$ at a fixed restart length
$k$.

\begin{definition}[Restarted GMRES, GMRES($k$)]
\label{def:restart}
Fix a restart length $k$. Run GMRES for at most $k$ steps. If the residual has not
met the tolerance by step $k$, \emph{restart}: discard the basis $\mat{V}_k$ and the
Hessenberg matrix entirely, take the current iterate $\vx_k$ as the new initial guess
$\vx_0$, recompute the residual $\vr_0=\vb-\mat{A}\vx_0$, and begin a fresh GMRES
cycle of up to $k$ steps. Repeat until convergence. Storage is capped at $O(kN)$ and
per-cycle work at $O(k^2N)$, both independent of the total iteration count.
\end{definition}

The basis is \emph{ephemeral}: it is born at the start of a cycle and discarded at
the restart boundary. This is exactly the lifetime story of \cref{ch:strata}---the
Krylov basis is per-solve, indeed per-\emph{cycle}, run-time state, allocated fresh
and thrown away, never part of any operator's frozen internals. Only the iterate
$\vx$ survives across a restart; everything else is rebuilt.

\begin{figure}[t]
\centering
\begin{tikzpicture}
  \begin{axis}[
      width=12.0cm, height=5.4cm,
      xlabel={total iterations},
      ylabel={basis size (vectors stored)},
      xmin=0, xmax=15, ymin=0, ymax=8,
      xtick={0,3,6,9,12,15}, ytick={0,2,4,6},
      tick label style={font=\scriptsize},
      label style={font=\small},
      axis lines=left, clip=false,
    ]
    % full GMRES: linear growth
    \addplot[simBlue, thick, mark=*, mark size=1.2pt] coordinates
      {(0,0)(1,1)(2,2)(3,3)(4,4)(5,5)(6,6)(7,7)};
    \node[font=\scriptsize, simBlue, anchor=west] at (axis cs:6.1,6.9)
      {full GMRES: $O(m)$};
    % restarted GMRES(3): sawtooth
    \addplot[simRust, very thick] coordinates
      {(0,0)(1,1)(2,2)(3,3)};
    \addplot[simRust, very thick] coordinates
      {(3,0)(4,1)(5,2)(6,3)};
    \addplot[simRust, very thick] coordinates
      {(6,0)(7,1)(8,2)(9,3)};
    \addplot[simRust, very thick] coordinates
      {(9,0)(10,1)(11,2)(12,3)};
    \addplot[simRust, very thick] coordinates
      {(12,0)(13,1)(14,2)(15,3)};
    % restart drop markers
    \draw[simRust, densely dashed] (axis cs:3,3) -- (axis cs:3,0);
    \draw[simRust, densely dashed] (axis cs:6,3) -- (axis cs:6,0);
    \draw[simRust, densely dashed] (axis cs:9,3) -- (axis cs:9,0);
    \draw[simRust, densely dashed] (axis cs:12,3) -- (axis cs:12,0);
    \node[font=\scriptsize, simRust, anchor=west] at (axis cs:9.3,3.5)
      {GMRES(3): sawtooth, capped at $k{=}3$};
  \end{axis}
\end{tikzpicture}
\caption[Memory: full GMRES versus the restart sawtooth]{Basis storage over the
course of a solve. Full GMRES (blue) accumulates one vector per step forever, so its
memory climbs linearly without bound. Restarted GMRES($k$) (rust, here $k=3$) caps
the basis at $k$ vectors: every $k$ steps it discards the basis and starts a fresh
cycle from the current iterate, producing the characteristic sawtooth. The peak
memory is bounded by $O(kN)$ no matter how many total iterations the solve takes. The
trade is convergence quality, charted in \cref{fig:convergence}.}
\label{fig:memory}
\end{figure}

\subsection{The trade-off}

Restarting is not free either; it trades memory for convergence speed.

\begin{pitfall}
\textbf{Restarting loses GMRES's optimality, and can stall.} Full GMRES minimizes
the residual over the \emph{entire} accumulated Krylov subspace---its convergence is
often superlinear, accelerating as the subspace captures more of the operator's
spectrum. The instant you restart, that accumulated subspace is thrown away: the new
cycle minimizes only over a \emph{fresh} $k$-dimensional space built from the current
residual, blind to everything the discarded basis had learned. The monotonic
\emph{decrease} of the residual still holds across restarts (each cycle starts from a
no-worse iterate), but the \emph{rate} can collapse---and for some operators
GMRES($k$) \emph{stagnates}, making negligible progress per cycle while full GMRES
would have converged. Choosing $k$ is a real engineering decision: too small and you
stall; too large and you run out of memory. There is no value of $k$ that is right
for every problem.
\end{pitfall}

\section{FGMRES: when the preconditioner changes every step}

We now pay off the promise of \cref{ch:operators}. There we met the
\emph{decision rule}: a variant bound once at construction is absorbed by a
constructed operator; a variant that changes \emph{per step} cannot be, and must
instead be threaded through the run-time state. We claimed that this single rule is
``exactly why \code{FGMRES} is a different algorithm from \code{GMRES}.'' Here is the
claim, made good.

\subsection{Where preconditioned GMRES applies $\mat{M}$}

A preconditioner $\mat{M}$ is an approximate inverse of $\mat{A}$ whose job is to
cluster the spectrum so the Krylov iteration converges faster (\cref{ch:precond}).
\emph{Right}-preconditioned GMRES solves $\mat{A}\mat{M}^{-1}\vu=\vb$ with
$\vx=\mat{M}^{-1}\vu$; concretely, each Arnoldi step preconditions the current basis
vector before the matvec: $\vz_j=\mat{M}^{-1}\vv_j$, then $\vw=\mat{A}\vz_j$. If
$\mat{M}$ is \emph{fixed}---chosen once at solve start and held constant---this is a
textbook constructed operator (\cref{wex:jacobi} in spirit): build it once, apply the
same \code{pc} at every step. The preconditioner is build-time configuration; it never
enters the threaded state. The left branch of the decision rule
(\cref{fig:decisionrule} of \cref{ch:operators}) handles it cleanly.

\subsection{What breaks when $\mat{M}$ varies}

Now allow the preconditioner to differ at each step: $\mat{M}_j$ at step $j$. This is
not exotic---it is exactly what happens when the preconditioner is \emph{itself an
inner iterative solve} (a few steps of an inner GMRES, a multigrid cycle run to a
loose, varying tolerance, a smoother whose accuracy adapts to the residual). Each
inner solve produces a slightly different linear operator, so $\vz_j=\mat{M}_j^{-1}
\vv_j$ uses a \emph{different} $\mat{M}_j$ for every $j$. There is no single $\mat{M}$
to construct.

Trace what fails. In right-preconditioned GMRES the iterate is recovered, at the end,
as $\vx=\vx_0+\mat{M}^{-1}\mat{V}_m\vy$---the basis $\mat{V}_m$ lives in the
\emph{preconditioned} space, and we map back by applying $\mat{M}^{-1}$ to the
correction. This reconstruction \emph{assumes a single fixed $\mat{M}$}: it applies
one $\mat{M}^{-1}$ to the whole combination $\mat{V}_m\vy$. But with per-step
$\mat{M}_j$, the $j$-th basis vector was preconditioned by $\mat{M}_j$, the next by
$\mat{M}_{j+1}$, and there is no single $\mat{M}^{-1}$ that undoes the mixture. The
operator that produced $\vz_j$ is \emph{gone} by the next step---it was an ephemeral
inner solve, not a stored value---so $\vz_j$ cannot be recomputed on demand. Plain
GMRES, which never stored the $\vz_j$, simply cannot form the right answer.

\begin{pitfall}
\textbf{Driving plain GMRES with a per-step preconditioner silently corrupts the
solve.} If you feed a varying $\mat{M}_j$ to an algorithm that recovers its iterate
as $\vx_0+\mat{M}^{-1}\mat{V}_m\vy$ under the assumption of a single fixed $\mat{M}$,
the code \emph{runs}---no crash, no error---and even reports a small ``residual''
$|s_{m+1}|$ from its running QR. But that residual is the residual of a least-squares
problem the iterate no longer corresponds to: the reconstruction used the wrong
inverse. The solve converges to the wrong answer, quietly. This is the silent partial
absorption of \cref{ch:operators} in its most dangerous form---a tidy equation hiding
an operation chain that has secretly forked. The repair is not a patch to GMRES; it is
a different algorithm that \emph{remembers} the preconditioned vectors. That algorithm
is FGMRES.
\end{pitfall}

\subsection{The repair: keep a second basis}

FGMRES---flexible GMRES---makes the only move the decision rule allows: it routes the
per-step variant into the threaded state by \emph{remembering} every preconditioned
vector. Alongside the Arnoldi basis $\mat{V}_m$ (the orthonormalized, Arnoldi-built
vectors), FGMRES carries a \emph{second} basis $\mat{Z}_m=[\vz_1\ \cdots\ \vz_m]$
holding the preconditioned vectors $\vz_j=\mat{M}_j^{-1}\vv_j$ exactly as they were
produced, each by its own ephemeral $\mat{M}_j$. The Arnoldi process and the running
QR are \emph{identical} to GMRES---byte for byte, the same orthogonalization, the same
Givens stream, the same free residual. Only two things change: at each step we save
$\vz_j$ into $\mat{Z}_m$, and at the end we reconstruct the iterate against
$\mat{Z}_m$ instead of against $\mat{V}_m$.

\begin{keyidea}
\textbf{The $\mat{Z}$-versus-$\mat{V}$ distinction.} GMRES recovers its iterate from
the Arnoldi basis, $\vx_m=\vx_0+\mat{V}_m\vy_m$. FGMRES recovers it from the stored
\emph{preconditioned} basis,
\[
  \vx_m \;=\; \vx_0 + \mat{Z}_m\,\vy_m,
  \qquad \mat{Z}_m=[\mat{M}_1^{-1}\vv_1\ \cdots\ \mat{M}_m^{-1}\vv_m].
\]
Because each $\vz_j=\mat{M}_j^{-1}\vv_j$ already carries the action of its own
per-step preconditioner, the reconstruction $\mat{Z}_m\vy_m$ never needs to apply a
single ``$\mat{M}^{-1}$'' to the combination---the per-step inverses are baked into the
columns of $\mat{Z}$. The flexibility is paid for by exactly one extra length-$N$ basis
in the threaded state. This is the ``\code{FgmresState} adds a $\mat{Z}$ field''
prediction of \cref{wex:flexible} (\cref{ch:operators}), realized.
\end{keyidea}

\begin{figure}[t]
\centering
\begin{tikzpicture}[scale=0.95, node distance=6mm]
  % --- the per-step pipeline, drawn once ---
  \node[font=\bfseries\small, simInk] at (0,2.9) {FGMRES threads two bases};
  % v_j
  \node[draw=simBlue, fill=simBgBlue, thick, rounded corners=2pt,
        minimum width=15mm, minimum height=8mm, font=\footnotesize] (vj) at (-4.4,1.4)
    {$\vv_j$};
  % the per-step preconditioner (ephemeral, dashed)
  \node[extern, minimum width=20mm, minimum height=9mm, right=8mm of vj] (Mj)
    {$\mat{M}_j^{-1}$ \scriptsize(per step)};
  % z_j
  \node[draw=simRust, fill=simBgRust, thick, rounded corners=2pt,
        minimum width=15mm, minimum height=8mm, font=\footnotesize, right=8mm of Mj] (zj)
    {$\vz_j$};
  % A
  \node[draw=simTeal, fill=simBgGray, thick, rounded corners=2pt,
        minimum width=12mm, minimum height=8mm, font=\footnotesize, right=8mm of zj] (A)
    {$\mat{A}$};
  % w
  \node[draw=simBlue, fill=simBgBlue, thick, rounded corners=2pt,
        minimum width=15mm, minimum height=8mm, font=\footnotesize, right=8mm of A] (w)
    {$\vw\!\to\!\vv_{j+1}$};
  \draw[flow] (vj) -- (Mj);
  \draw[flow] (Mj) -- (zj);
  \draw[flow] (zj) -- (A);
  \draw[flow] (A) -- node[above,font=\scriptsize]{matvec} (w);
  % the two bases being accumulated
  \node[draw=simBlue, fill=simBgBlue, semithick, rounded corners=2pt,
        minimum width=70mm, minimum height=7mm, font=\scriptsize, align=center] (Vbasis)
    at (0,-0.4) {Arnoldi basis $\mat{V}_m=[\vv_1\,\cdots\,\vv_m]$
                 \quad\itshape (orthonormal, built by Arnoldi)};
  \node[draw=simRust, fill=simBgRust, semithick, rounded corners=2pt,
        minimum width=70mm, minimum height=7mm, font=\scriptsize, align=center] (Zbasis)
    at (0,-1.5) {preconditioned basis $\mat{Z}_m=[\vz_1\,\cdots\,\vz_m]$
                 \quad\itshape (the per-step $\mat{M}_j^{-1}\vv_j$, remembered)};
  \draw[refedge] (vj.south) to[out=-90,in=90] (Vbasis.north -| -3.3,-0.4);
  \draw[refedge] (zj.south) to[out=-90,in=90] (Zbasis.north -| 0.5,-1.5);
  % the reconstruction
  \node[font=\scriptsize, simRust, align=center] at (0,-2.5)
    {iterate recovered from $\mat{Z}$, not $\mat{V}$:\quad
     $\vx_m=\vx_0+\mat{Z}_m\vy_m$};
\end{tikzpicture}
\caption[FGMRES's two bases]{Why FGMRES carries a second basis. Each step
preconditions the Arnoldi vector $\vv_j$ by a \emph{per-step} operator $\mat{M}_j^{-1}$
(dashed: ephemeral, gone by the next step---perhaps an inner iterative solve), giving
$\vz_j$, which is then hit with $\mat{A}$ to continue the Arnoldi process. The
algorithm accumulates \emph{both} bases: $\mat{V}_m$ (the orthonormal Arnoldi vectors,
blue) drives orthogonalization and the running QR exactly as in plain GMRES; but the
iterate is reconstructed from $\mat{Z}_m$ (the remembered preconditioned vectors,
rust), because each $\vz_j$ already carries its own $\mat{M}_j^{-1}$ and no single
$\mat{M}^{-1}$ could undo the mixture. The per-step variant has become a field of the
threaded state---precisely the right branch of the decision rule
(\cref{fig:decisionrule}, \cref{ch:operators}).}
\label{fig:fgmres}
\end{figure}

\Cref{fig:fgmres} draws the doubled basis. The structural reading is exactly the one
\cref{ch:operators} set up: the variant ``which preconditioner this step'' has a
per-step lifetime, so by the decision rule it cannot live in a constructed operator;
it lives in the carry, as the extra basis $\mat{Z}$. FGMRES is GMRES with one field
added to its state schema---and that field is forced, not chosen, by the lifetime of
the varying preconditioner. The least-squares machinery does not even notice: in the
source, the FGMRES running-QR stream is line-for-line identical to the GMRES one; the
\emph{only} difference is the back-solve reconstructs the correction against $\mat{Z}$
rather than $\mat{V}$.

\begin{workedexample}{GMRES versus restarted GMRES(2): convergence on one system}{restart}
We compare full GMRES with restarted GMRES(2) on a $4\times4$ nonsymmetric system,
watching the residual norms (read for free from $|s_{m+1}|$) at each step. The system
is mildly nonnormal, so restarting visibly hurts.

\textbf{Full GMRES.} Storing the whole basis, full GMRES drives the residual down
monotonically and, on a $4\times4$ system, reaches the exact solution by step $4$ (in
exact arithmetic, GMRES converges in at most $N$ steps because $\mathcal{K}_N$ spans
the whole space). A representative residual history:
\[
  \norm{\vr_0}=1,\quad \norm{\vr_1}\approx0.42,\quad \norm{\vr_2}\approx0.15,\quad
  \norm{\vr_3}\approx0.03,\quad \norm{\vr_4}\approx10^{-15}.
\]

\textbf{Restarted GMRES(2).} Capped at $k=2$, GMRES(2) discards its basis after every
second step. The first cycle matches full GMRES through step $2$
($\norm{\vr_2}\approx0.15$), but then it \emph{restarts}: the rich 2-dimensional
subspace is thrown away, and cycle 2 rebuilds a fresh 2-dimensional space from the
step-2 residual---blind to what cycle 1 learned. Progress per cycle slows:
\[
  0.15 \;\to\; 0.066 \;\to\; 0.030 \;\to\; 0.014 \;\to\; 0.0065 \;\to\;\cdots
\]
roughly halving each cycle---a \emph{linear} rate, where full GMRES was racing to
machine precision. \Cref{fig:convergence} plots both. GMRES(2) still converges (the
residual is monotone), but it takes many more iterations to reach the same tolerance:
the price of bounding memory at $k=2$ rather than letting the basis grow to $N=4$.

\textbf{The reading.} Full GMRES exhausts the Krylov space and hits the exact answer
at step $N$; restarted GMRES(2) trades that finite-termination guarantee, and the
superlinear tail, for a flat $O(kN)$ memory ceiling. On this small system the gap is
stark; on a million-unknown system where full GMRES is impossible, GMRES($k$)'s slow
linear convergence is simply the cost of being able to run at all.
\end{workedexample}

\begin{figure}[t]
\centering
\begin{tikzpicture}
  \begin{axis}[
      width=12.2cm, height=6.2cm,
      xlabel={iteration $m$}, ylabel={residual norm $\norm{\vr_m}$},
      ymode=log, ymin=1e-16, ymax=2, xmin=0, xmax=10,
      xtick={0,2,4,6,8,10},
      ytick={1e0,1e-3,1e-6,1e-9,1e-12,1e-15},
      tick label style={font=\scriptsize}, label style={font=\small},
      legend style={font=\scriptsize, at={(0.98,0.98)}, anchor=north east},
      grid=both, grid style={simGray!18},
      axis lines=left, clip=false,
    ]
    % full GMRES: superlinear, hits machine precision at m=4
    \addplot[simBlue, very thick, mark=*, mark size=1.4pt] coordinates
      {(0,1)(1,0.42)(2,0.15)(3,0.03)(4,1e-15)};
    \addlegendentry{full GMRES}
    % restarted GMRES(2): linear, restarts every 2 steps (dashed verticals)
    \addplot[simRust, very thick, mark=square*, mark size=1.3pt] coordinates
      {(0,1)(1,0.42)(2,0.15)(3,0.066)(4,0.030)(5,0.014)(6,0.0065)(7,0.0030)(8,0.0014)(9,0.00065)(10,0.00030)};
    \addlegendentry{restarted GMRES(2)}
    % restart boundaries
    \draw[simRust, densely dotted] (axis cs:2,1e-16) -- (axis cs:2,2);
    \draw[simRust, densely dotted] (axis cs:4,1e-16) -- (axis cs:4,2);
    \draw[simRust, densely dotted] (axis cs:6,1e-16) -- (axis cs:6,2);
    \draw[simRust, densely dotted] (axis cs:8,1e-16) -- (axis cs:8,2);
    \node[font=\scriptsize, simRust, anchor=south] at (axis cs:4,1.1) {restarts};
  \end{axis}
\end{tikzpicture}
\caption[Convergence: full GMRES versus GMRES(2)]{Residual histories for
\cref{wex:restart} (log scale). Full GMRES (blue) converges superlinearly and reaches
machine precision at step $4=N$, exhausting the Krylov space. Restarted GMRES(2)
(rust) matches it for the first two steps, then restarts (dotted verticals) every two
steps; each restart discards the accumulated subspace, collapsing the convergence to a
slow linear rate. Both residual curves are monotone non-increasing---restart never
\emph{increases} the residual---but GMRES(2) needs many more iterations to reach the
same tolerance. The flat $O(kN)$ memory ceiling is paid for in convergence speed.}
\label{fig:convergence}
\end{figure}

\section{The variant axes: preconditioner side and the absorption view}

GMRES carries several orthogonal axes of variation, and they line up precisely with
the variant-absorption discipline of \cref{ch:operators}. We collect them here.

\subsection{Left, right, and no preconditioning}

A preconditioner can be applied on either side of $\mat{A}$, or not at all
(\cref{fig:precondside}):
\begin{description}[style=nextline,leftmargin=1.5em]
  \item[None.] Solve $\mat{A}\vx=\vb$ directly; the Krylov space is built from
  $\mat{A}$ and $\vr_0$. Used when $\mat{A}$ is already well-conditioned.
  \item[Left.] Solve $\mat{M}^{-1}\mat{A}\vx=\mat{M}^{-1}\vb$. The matvec at each step
  is $\vw=\mat{M}^{-1}(\mat{A}\vv_j)$ (precondition \emph{after} the operator). The
  residual GMRES minimizes is the \emph{preconditioned} residual
  $\norm{\mat{M}^{-1}(\vb-\mat{A}\vx)}$---which is not the true residual, a subtlety
  when reading the convergence test.
  \item[Right.] Solve $\mat{A}\mat{M}^{-1}\vu=\vb$ with $\vx=\mat{M}^{-1}\vu$. The
  matvec is $\vw=\mat{A}(\mat{M}^{-1}\vv_j)$ (precondition \emph{before} the operator).
  The residual minimized is the \emph{true} residual $\norm{\vb-\mat{A}\vx}$, which is
  why right preconditioning is often preferred for GMRES---and why FGMRES, which
  preconditions before the matvec, is naturally a right-preconditioned scheme.
\end{description}

\begin{figure}[t]
\centering
\begin{tikzpicture}[node distance=5mm]
  % three little pipelines stacked
  \foreach \row/\lbl/\seq in {
      1.7/{none}/{$\vv_j$/$\mat{A}$/$\vw$},
      0.0/{left}/{$\vv_j$/$\mat{A}$/$\mat{M}^{-1}$/$\vw$},
      -1.7/{right}/{$\vv_j$/$\mat{M}^{-1}$/$\mat{A}$/$\vw$}}{
    \node[font=\footnotesize\bfseries, simInk, anchor=east] at (-4.3,\row) {\lbl:};
  }
  % NONE
  \node[draw=simBlue,fill=simBgBlue,thick,rounded corners=2pt,minimum width=11mm,minimum height=7mm,font=\footnotesize] (n1) at (-3.0,1.7) {$\vv_j$};
  \node[draw=simTeal,fill=simBgGray,thick,rounded corners=2pt,minimum width=11mm,minimum height=7mm,font=\footnotesize,right=7mm of n1] (n2) {$\mat{A}$};
  \node[draw=simBlue,fill=simBgBlue,thick,rounded corners=2pt,minimum width=11mm,minimum height=7mm,font=\footnotesize,right=7mm of n2] (n3) {$\vw$};
  \draw[flow] (n1)--(n2); \draw[flow] (n2)--(n3);
  % LEFT
  \node[draw=simBlue,fill=simBgBlue,thick,rounded corners=2pt,minimum width=11mm,minimum height=7mm,font=\footnotesize] (l1) at (-3.0,0.0) {$\vv_j$};
  \node[draw=simTeal,fill=simBgGray,thick,rounded corners=2pt,minimum width=11mm,minimum height=7mm,font=\footnotesize,right=7mm of l1] (l2) {$\mat{A}$};
  \node[extern,minimum width=12mm,minimum height=7mm,font=\footnotesize,right=7mm of l2] (l3) {$\mat{M}^{-1}$};
  \node[draw=simBlue,fill=simBgBlue,thick,rounded corners=2pt,minimum width=11mm,minimum height=7mm,font=\footnotesize,right=7mm of l3] (l4) {$\vw$};
  \draw[flow] (l1)--(l2); \draw[flow] (l2)--(l3); \draw[flow] (l3)--(l4);
  % RIGHT
  \node[draw=simBlue,fill=simBgBlue,thick,rounded corners=2pt,minimum width=11mm,minimum height=7mm,font=\footnotesize] (r1) at (-3.0,-1.7) {$\vv_j$};
  \node[extern,minimum width=12mm,minimum height=7mm,font=\footnotesize,right=7mm of r1] (r2) {$\mat{M}^{-1}$};
  \node[draw=simTeal,fill=simBgGray,thick,rounded corners=2pt,minimum width=11mm,minimum height=7mm,font=\footnotesize,right=7mm of r2] (r3) {$\mat{A}$};
  \node[draw=simBlue,fill=simBgBlue,thick,rounded corners=2pt,minimum width=11mm,minimum height=7mm,font=\footnotesize,right=7mm of r3] (r4) {$\vw$};
  \draw[flow] (r1)--(r2); \draw[flow] (r2)--(r3); \draw[flow] (r3)--(r4);
  % note
  \node[font=\scriptsize\itshape, simGray, align=left, anchor=west] at (3.6,1.7)
    {minimizes true residual};
  \node[font=\scriptsize\itshape, simGray, align=left, anchor=west] at (3.6,0.0)
    {minimizes $\mat{M}^{-1}$-residual};
  \node[font=\scriptsize\itshape, simGray, align=left, anchor=west] at (3.6,-1.7)
    {minimizes true residual};
\end{tikzpicture}
\caption[The preconditioner-side axis]{The three settings of the preconditioner-side
variant. \emph{None:} the matvec is just $\mat{A}\vv_j$. \emph{Left:} apply
$\mat{M}^{-1}$ \emph{after} $\mat{A}$, solving $\mat{M}^{-1}\mat{A}\vx=\mat{M}^{-1}\vb$;
the residual GMRES sees is the preconditioned one. \emph{Right:} apply $\mat{M}^{-1}$
\emph{before} $\mat{A}$, solving $\mat{A}\mat{M}^{-1}\vu=\vb$; the residual stays the
true one. All three are the \emph{same} Arnoldi-plus-running-QR loop with the matvec's
operand chain rebound---a textbook variant absorbed by a constructed operator
(\cref{ch:operators}): the side is read once, at construction, and the per-step loop
never branches on it.}
\label{fig:precondside}
\end{figure}

\subsection{The absorption view}

These axes are the running example of \cref{ch:operators}'s discipline. The
preconditioner side, the restart length, the orthogonalization method, even the
choice of CG/GMRES/FGMRES itself are all read \emph{once}, at the factory that
constructs the solver (\code{ConfigureKrylovSolver}), and absorbed into a single
uniformly typed solver-operator. After construction the per-step loop is
variant-free: it asks the operator for \code{apply(A, v)} and the preconditioner for
\code{apply(pc, r)}, and never re-inspects which side, which method, which restart it
got. The fixed-versus-flexible axis is the \emph{one} axis that cannot be absorbed
this way---because, by the decision rule, a per-step preconditioner has run-time
lifetime---and that is exactly why FGMRES is a distinct algorithm with a distinct
state schema rather than a configuration flag on GMRES. Everything else is a binding;
flexibility alone is a structural fork, honestly disclosed by the extra $\mat{Z}$
basis.

\subsection{Convergence intuition: no clean condition-number bound}

CG (\cref{ch:cg}) came with a clean convergence bound: the error after $m$ steps
shrinks like $\bigl(\tfrac{\sqrt\kappa-1}{\sqrt\kappa+1}\bigr)^m$, governed entirely
by the condition number $\kappa$. GMRES has \emph{no such bound}, and the reason is
instructive. For a symmetric operator the spectrum lives on the real line and the
condition number captures everything; for a nonsymmetric operator the eigenvalues
scatter into the complex plane, and---worse---if $\mat{A}$ is far from normal, the
eigenvalues can lie anywhere and \emph{still} mislead, because the relevant object is
not the spectrum but the \emph{field of values} (the set $\{\inner{\mat{A}\vx}{\vx}:
\norm{\vx}=1\}$, a region of the complex plane). GMRES converges fast when the field
of values is a tidy cluster bounded away from the origin, and can stagnate when it
wraps around the origin---regardless of the condition number. The practical
consequence: GMRES convergence is judged empirically, by watching the free residual
$|s_{m+1}|$ fall, and is engineered by \emph{preconditioning}---reshaping the field of
values with $\mat{M}$ until the cluster is favorable. The whole point of
\cref{ch:precond,ch:multigrid} is to build the $\mat{M}$ that makes GMRES converge.

\summarysection
When $\mat{A}$ is nonsymmetric, the short three-term recurrence that makes CG cheap
simply does not exist (Faber--Manteuffel): optimality and bounded storage cannot
coexist. \emph{GMRES} keeps optimality and stores the whole basis. It is built on the
\emph{minimum-residual principle}---at step $m$, return the $\vx_m\in\vx_0+\mathcal{K}_m$
minimizing $\norm{\vb-\mat{A}\vx_m}$---which makes the residual norms monotone
non-increasing. The \emph{Arnoldi process} (Gram--Schmidt on the Krylov sequence,
\cref{ch:orthog}) builds an orthonormal basis $\mat{V}_{m+1}$ and the small
upper-Hessenberg $\bar{\mat{H}}_m$ satisfying the Arnoldi relation
$\mat{A}\mat{V}_m=\mat{V}_{m+1}\bar{\mat{H}}_m$; orthonormality then collapses the
$N$-dimensional minimization to the $(m{+}1)$-dimensional least-squares problem
$\min_{\vy}\norm{\beta\vec e_1-\bar{\mat{H}}_m\vy}$. A \emph{running QR by Givens
rotations} (\cref{ch:givens}) solves that small problem incrementally and---the central
economy---exposes the residual norm \emph{for free} as $|s_{m+1}|$ at every step,
without forming the iterate. Because storage grows as $O(mN)$ and work as $O(m^2N)$,
\emph{restarted GMRES($k$)} caps memory at $O(kN)$ by discarding the (ephemeral,
\cref{ch:strata}) basis every $k$ steps---at the cost of losing superlinear
convergence and risking stagnation. Finally, \emph{FGMRES} pays off \cref{ch:operators}'s
decision rule: when the preconditioner changes per step ($\mat{M}_j$, e.g.\ an inner
iterative solve), no constructed operator can absorb it, so the per-step variant is
threaded into the state as a second basis $\mat{Z}_m$ of remembered preconditioned
vectors, and the iterate is recovered as $\vx_m=\vx_0+\mat{Z}_m\vy_m$ rather than
$\vx_0+\mat{V}_m\vy_m$. Preconditioner side (left/right/none) is an absorbed variant; the
fixed-versus-flexible axis is the one structural fork, disclosed honestly by the extra
basis. GMRES carries no clean condition-number bound: its convergence is governed by the
field of values and engineered by preconditioning.

\furtherreading
The definitive treatment of GMRES, FGMRES, the Arnoldi process, restarting, and the
preconditioning variants is Saad~\citep{saad2003}, whose chapters on Krylov methods for
nonsymmetric systems and on flexible preconditioning are the direct source for this
chapter's algorithm and its decision-rule reading. For the numerical-analysis
perspective---the Arnoldi relation as an orthogonal reduction, the role of
nonnormality and the field of values in convergence, and the contrast with the clean
symmetric theory---Trefethen and Bau~\citep{trefethenbau1997} (the lectures on
Arnoldi, GMRES, and iterative-method convergence) is the ideal companion. The
running-QR/Givens mechanics this chapter used as a black box are developed in full in
\cref{ch:givens}; the orthogonalization inside Arnoldi in \cref{ch:orthog}; and the
preconditioners $\mat{M}$ whose construction makes GMRES converge in
\cref{ch:precond,ch:multigrid}.

\exercisessection
\begin{enumerate}[nosep]
  \item \textbf{Why no short recurrence.} State, in your own words, what the
  Faber--Manteuffel theorem rules out, and explain why CG's three-term recurrence is a
  special consequence of symmetry. What \emph{two} desirable properties does GMRES
  refuse to give up simultaneously, and which one does restarting trade away?

  \item \textbf{Monotonicity.} Prove that the GMRES residual norms satisfy
  $\norm{\vr_{m+1}}\le\norm{\vr_m}$ using only the minimum-residual principle and the
  nesting $\mathcal{K}_m\subseteq\mathcal{K}_{m+1}$. Why does this argument
  \emph{not} establish strict decrease, and when can the residual stall at a plateau?

  \item \textbf{A third Arnoldi step.} Continue \cref{wex:arnoldi} with one more step:
  compute $\mat{A}\vv_3$, the coefficients $h_{13},h_{23},h_{33}$, and the new
  sub-diagonal $h_{43}$. What is special about $h_{43}$ for this particular $3\times3$
  operator, and what does its value tell you about where GMRES terminates exactly?

  \item \textbf{The reduction, by hand.} Starting from $\vr_0=\beta\vv_1$ and the
  Arnoldi relation, rederive the identity $\vb-\mat{A}(\vx_0+\mat{V}_m\vy)=
  \mat{V}_{m+1}(\beta\vec e_1-\bar{\mat{H}}_m\vy)$, naming the single property of
  $\mat{V}_{m+1}$ that lets you drop it from the norm. Which step would fail if the
  Arnoldi basis were merely linearly independent rather than orthonormal?

  \item \textbf{The free residual.} In \cref{wex:leastsq} the running QR reported
  $\norm{\vr_2}=|s_3|=1/\sqrt{69}$ \emph{without} forming $\vx_2$. Explain precisely
  why $|s_{m+1}|$ equals the least-squares residual norm, citing the isometry property
  of Givens rotations. Then state what GMRES would have to compute instead if it lacked
  this readout, and estimate the cost difference per step.

  \item \textbf{Sizing a restart.} A solve has $N=10^6$ unknowns; each vector is $8$~MB.
  You have a $4$~GB memory budget for the Krylov basis. (i) What is the largest restart
  length $k$ you can afford? (ii) If full GMRES would converge in $40$ steps but
  GMRES($k$) for your $k$ needs $300$, has restarting helped or hurt, and in what sense?
  (iii) Name one operator property under which GMRES($k$) might \emph{stagnate} entirely.

  \item \textbf{The mutating-preconditioner trap, revisited.} A colleague proposes to
  run \emph{plain} GMRES with an inner-iterative preconditioner whose accuracy varies
  step to step, ``since GMRES only ever asks \code{apply(pc, v)} and won't notice.''
  Using the $\mat{Z}$-versus-$\mat{V}$ distinction and the decision rule of
  \cref{ch:operators}, explain exactly what goes wrong, why no error is raised, and why
  the reported $|s_{m+1}|$ is misleading. What is the minimal change that fixes it, and
  what does it cost in storage?

  \item \textbf{$\mat{V}$ versus $\mat{Z}$.} For \emph{fixed} right-preconditioned GMRES,
  the iterate is $\vx_0+\mat{M}^{-1}\mat{V}_m\vy$; for FGMRES it is $\vx_0+\mat{Z}_m\vy$.
  Show that when $\mat{M}_j=\mat{M}$ for all $j$ (the preconditioner happens to be
  constant), these two formulas \emph{coincide}, so FGMRES reduces to fixed
  right-preconditioned GMRES. What, then, is the only cost of always using FGMRES, and
  when is paying it worthwhile?

  \item \textbf{Side and the residual.} Left and right preconditioning minimize
  \emph{different} residuals (\cref{fig:precondside}). For a convergence test with
  tolerance $\varepsilon$ on the \emph{true} residual $\norm{\vb-\mat{A}\vx}$, which
  side lets you read the test directly from the free $|s_{m+1}|$, and which forces a
  correction? Explain why this is one practical reason GMRES implementations often
  default to right preconditioning.

  \item \textbf{Absorbed or forked?} Classify each axis as either \emph{absorbed} (read
  once at construction, per-step loop variant-free) or a \emph{structural fork} (forces
  a different operation chain or state schema), and justify with the decision rule:
  (a) preconditioner side $\in\{\text{none},\text{left},\text{right}\}$;
  (b) restart length $k$; (c) orthogonalization method (MGS vs CGS); (d)
  fixed vs per-step preconditioner. For the one you classify as a fork, name the extra
  state it forces into the carry.
\end{enumerate}
